%!TEX program = xelatex
%!TEX encoding = UTF-8
%Definições de Ano, Semestre, número e data da avaliação
\def\numeroteste{1}
\def\modulo{2}
\def\dataavaliacao{11/10/2023}

\documentclass[12pt]{exam}

\usepackage[singlelinecheck=false]{caption}
\usepackage{textpos}
\usepackage{amssymb}
\usepackage{amsmath,amsfonts,amsthm,amstext}
\usepackage[justification=centering]{caption}
\usepackage[brazil]{babel}
% \usepackage[latin1]{inputenc}
\usepackage{graphicx}
\graphicspath{{../../figuras_geogebra/}{../../Pictures/}}
\usepackage{enumitem}
\usepackage{multicol}
\usepackage{answers}
\usepackage{qrcode}

\usepackage{pifont}
\newcommand{\cmark}{\ding{228}}%

% requires version 0.3 of the package
\usepackage[customcolors,beamer]{hf-tikz}

\tikzset{style green/.style={
        set fill color=green!50!lime!60,
        set border color=white,
    },
    style cyan/.style={
        set fill color=cyan!90!blue!60,
        set border color=white,
    },
    style orange/.style={
        set fill color=orange!80!red!60,
        set border color=white,
    },
    hor/.style={
        above left offset={-0.15,0.31},
        below right offset={0.15,-0.125},
        #1
    },
    ver/.style={
        above left offset={-0.1,0.3},
        below right offset={0.15,-0.15},
        #1
    }
}

\usepackage{tikz,ifthen,tikz-3dplot}
\usepackage{siunitx}
\usetikzlibrary{circuits.ee.IEC}
\usetikzlibrary{lindenmayersystems}
\usetikzlibrary[shadings]
\usepackage{pgfplots}
\pgfplotsset{compat=1.15}
\usepackage{mathrsfs}
\usetikzlibrary{arrows}
\usepackage{tcolorbox}
\newcounter{saveenumi}
\newcommand{\seti}{\setcounter{saveenumi}{\value{enumi}}}
\newcommand{\conti}{\setcounter{enumi}{\value{saveenumi}}}

\resetcounteronoverlays{saveenumi}


\renewcommand{\qedsymbol}{$\blacksquare$}

%\setbeamertemplate{itemize items}[ball]

%\setbeamertemplate{itemize items}[default]
%-------------amatrix
% Augmented matrix.  Usage (note the argument does not count the aug col):
% \begin{amatrix}{2}
    %   1  2  3 \\  4  5  6
    % \end{amatrix}
\newenvironment{amatrix}[1]{%
    \left[\begin{array}{@{}*{#1}{r}|r@{}}
    }{%
    \end{array}\right]
}

\def\ano{2023}
\def\semestre{2}
\def\turma{11}
\newcommand{\nsub}{\varsubsetneq}
\newcommand{\vaz}{\emptyset}
\newcommand{\dom}{{\rm dom\,}}
\newcommand{\im}{{\rm Im\,}}
\newcommand{\sub}{\subseteq}
\newcommand{\n}{\mathbb{N}}
\newcommand{\z}{\mathbb{Z}}
\newcommand{\rac}{\mathbb{Q}}
\newcommand{\real}{\mathbb{R}}
\newcommand{\complex}{\mathbb{C}}
\newcommand{\cp}[1]{\mathbb{#1}}
\newcommand{\ch}{\mbox{\textrm{car\,}}\nobreak}
\newcommand{\p}[1]{\mbox{\textrm{posto(#1)}}}
\newcommand{\vesp}[1]{\vspace{ #1  cm}}
\newcommand{\compcent}[1]{\vcenter{\hbox{$#1\circ$}}}
\newcommand{\comp}{\mathbin{\mathchoice
{\compcent\scriptstyle}{\compcent\scriptstyle}
{\compcent\scriptscriptstyle}{\compcent\scriptscriptstyle}}}
\newcommand\T{\rule{0pt}{2.6ex}}
\newtheorem{definicao}{Defini\c{c}\~ao}[section]
\newtheorem{definicoes}{Defini\c{c}\~oes}[section]
\newtheorem{exemplo}{Exemplo}[section]
\newtheorem{exemplos}{Exemplos}[section]
\newtheorem{exercicio}{Exerc{\'\i}cio}
\newtheorem{observacao}{Observa{\c c}{\~a}o:}[section]
\newtheorem{observacoes}{Observa{\c c}{\~o}es:}[section]
\newtheorem*{solucao}{Solu{\c c}{\~a}o:}
\newtheorem{proposicao}{Proposi\c{c}\~ao}
\newtheorem{lema}{Lema}
\newtheorem{teorema}{Teorema}
\newtheorem{corolario}{Corol\'ario}
\newenvironment{prova}[1][Prova]{\noindent\textbf{#1:} }{\qedsymbol}%{\ \rule{0.5em}{0.5em}}

\begin{document}

    \begin{figure}[h]
        \begin{minipage}[c]{1.7cm}
            \includegraphics[width=1.7cm]{unb.pdf}
        \end{minipage}
        \hspace{0pt}
        \begin{minipage}[c]{4in}
            {Universidade de Brasília} \\
            {Departamento de Matemática}
        \end{minipage}
    \end{figure}
    \vesp{-1}
    \hrule
    \begin{center}
        {\Large\bf \disciplina\ - Turma \turma}  \\
        \vesp{0.2} {\large\bf Teste \numeroteste\ - Módulo\ \modulo\ -\ \dataavaliacao}
    \end{center}

    \noindent{\bf \underline{Nome: \hspace{9.8cm} Mat.:\qquad \
            \hspace{2.1cm}}}\\
    \vspace*{.01cm}

    \noindent{\bf \underline{Nome: \hspace{9.8cm} Mat.:\qquad \
            \hspace{2.1cm}}}
    \vspace{.4cm}

    \questao{} Seja
    \[
        V = M_2(\real) = \left\{\begin{bmatrix}a_{11} & a_{12}\\a_{21} & a_{22} \end{bmatrix} \mid a_{11}, a_{12}, a_{21}, a_{22} \in \real\right\}.
    \]
    Considere em $V$ as seguintes operações
    \begin{itemize}
        \item $A + B = \begin{bmatrix}a_{11} & a_{12}\\a_{21} & a_{22} \end{bmatrix} + \begin{bmatrix}b_{11} & b_{12}\\b_{21} & b_{22} \end{bmatrix} = \begin{bmatrix}a_{11} + b_{11 }& a_{12} + b_{12}\\a_{21} + b_{21} & a_{22} + b_{22}\end{bmatrix}$
        \item $\alpha \cdot A = \alpha \cdot \begin{bmatrix}a_{11} & a_{12}\\a_{21} & a_{22} \end{bmatrix} = \begin{bmatrix}\alpha a_{11} & \alpha a_{12}\\\alpha a_{21} & \alpha a_{22} \end{bmatrix}$
    \end{itemize}
    para todos $A = \begin{bmatrix}a_{11} & a_{12}\\a_{21} & a_{22} \end{bmatrix}$, $B = \begin{bmatrix}b_{11} & b_{12}\\b_{21} & b_{22} \end{bmatrix} \in V$ e todo $\alpha \in \cp{K} = \real$.
    Mostre que com essas operações $V = M_2(\real)$ é um espaço vetorial sobre $\cp{K} = \real$.

    \solucao

    De fato,
    \begin{itemize}
        \item[A1)] Sejam
        \[
        A = \begin{bmatrix} a_{11} & a_{12}\\a_{21} & a_{22}\end{bmatrix},
        B = \begin{bmatrix} b_{11} & b_{12}\\b_{21} & b_{22}\end{bmatrix}.
        \]
        Temos:
        \begin{align*}
            A + B & = \begin{bmatrix} a_{11} & a_{12}\\a_{21} & a_{22}\end{bmatrix} + \begin{bmatrix} b_{11} & b_{12}\\b_{21} & b_{22}\end{bmatrix}
            \\ &= \begin{bmatrix} a_{11} + b_{11} & a_{12} + b_{12}\\a_{21} + b_{21} & a_{22} + b_{22}\end{bmatrix}
            \\ &= \begin{bmatrix} b_{11} + a_{11} & b_{12} + a_{12}\\b_{21} + a_{21} & b_{22} + a_{22}\end{bmatrix}
            \\ &= \begin{bmatrix} b_{11} & b_{12}\\b_{21} & b_{22}\end{bmatrix} + \begin{bmatrix} a_{11} & a_{12}\\a_{21} & a_{22}\end{bmatrix}
            \\ &= B + A
        \end{align*}

        \item[A2)] Sejam
        \[
        A = \begin{bmatrix} a_{11} & a_{12}\\a_{21} & a_{22}\end{bmatrix},
        B = \begin{bmatrix} b_{11} & b_{12}\\b_{21} & b_{22}\end{bmatrix},
        C = \begin{bmatrix} c_{11} & c_{12}\\c_{21} & c_{22}\end{bmatrix}.
        \]
        Temos:
        \begin{align*}
            (A + B) + C & = \left(\begin{bmatrix} a_{11} & a_{12}\\a_{21} & a_{22}\end{bmatrix} +
            \begin{bmatrix} b_{11} & b_{12}\\b_{21} & b_{22}\end{bmatrix}\right) +
            \begin{bmatrix} c_{11} & c_{12}\\c_{21} & c_{22}\end{bmatrix}
            \\ &= \begin{bmatrix} a_{11} + b_{11} & a_{12} + b_{12}\\a_{21} + b_{21} & a_{22} + b_{22}\end{bmatrix} +
            \begin{bmatrix} c_{11} & c_{12}\\c_{21} & c_{22}\end{bmatrix}
            \\ &= \begin{bmatrix} (a_{11} + b_{11}) + c_{11} & (a_{12} + b_{12}) + c_{12}\\(a_{21} + b_{21}) + c_{21} & (a_{22} + b_{22}) + c_{22}\end{bmatrix}
            \\ &= \begin{bmatrix} a_{11} + (b_{11} + c_{11}) & a_{12} + (b_{12} + c_{12})\\a_{21} + (b_{21} + c_{21}) & a_{22} + (b_{22} + c_{22})\end{bmatrix}
            \\ &= \begin{bmatrix} a_{11} & a_{12}\\a_{21} & a_{22}\end{bmatrix} + \begin{bmatrix} b_{11} + c_{11} & b_{12} + c_{12}\\b_{21} + c_{21} & b_{22} + c_{22}\end{bmatrix}
            \\ &= \begin{bmatrix} a_{11} & a_{12}\\a_{21} & a_{22}\end{bmatrix} +
            \left(\begin{bmatrix} b_{11} & b_{12}\\b_{21} & b_{22}\end{bmatrix} +
            \begin{bmatrix} c_{11} & c_{12}\\c_{21} & c_{22}\end{bmatrix}\right)
            \\ &= A + (B + C).
        \end{align*}

        \item[A3)] Tome
        \[
        0_V = \begin{bmatrix}0 & 0\\0 & 0\end{bmatrix}.
        \]
        Para toda matriz
        \[
        A = \begin{bmatrix} a_{11} & a_{12}\\a_{21} & a_{22}\end{bmatrix},
        \]
        temos
        \begin{align*}
            A + 0_V & = \begin{bmatrix} a_{11} & a_{12}\\a_{21} & a_{22}\end{bmatrix} +
            \begin{bmatrix}0 & 0\\0 & 0\end{bmatrix}
            = \begin{bmatrix}a_{11} + 0 & a_{12} + 0\\a_{21} + 0  & a_{22} + 0\end{bmatrix}
            \\ &= \begin{bmatrix} a_{11} & a_{12}\\a_{21} & a_{22}\end{bmatrix} = A.
        \end{align*}
        Portanto a matriz
        \[
        0_V = \begin{bmatrix}0 & 0\\0 & 0\end{bmatrix}
        \]
        é o vetor nulo.

        \item[A4)] Dada
        \[
        A = \begin{bmatrix} a_{11} & a_{12}\\a_{21} & a_{22}\end{bmatrix} \in \cp{M}_{2 \times 2}(\real),
        \]
        tome
        \[
        -A = \begin{bmatrix} -a_{11} & -a_{12}\\-a_{21} & -a_{22}\end{bmatrix} \in \cp{M}_{2 \times 2}(\real).
        \]
        Assim
        \begin{align*}
            A + (-A) & = \begin{bmatrix} a_{11} & a_{12}\\a_{21} & a_{22}\end{bmatrix} +
            \begin{bmatrix} -a_{11} & -a_{12}\\-a_{21} & -a_{22}\end{bmatrix}
            \\ &= \begin{bmatrix} a_{11} - a_{11} & a_{12} - a_{12}\\a_{21} - a_{21} & a_{22} - a_{22}\end{bmatrix}
            \\ &= \begin{bmatrix}0 & 0\\0 & 0\end{bmatrix} = 0_V.
        \end{align*}
        Logo a matriz
        \[
        -A = \begin{bmatrix} -a_{11} & -a_{12}\\-a_{21} & -a_{22}\end{bmatrix}
        \]
        é o vetor oposto.

        \item[M1)] Sejam $\alpha$, $\beta \in \real$ e
        \[
        A = \begin{bmatrix} a_{11} & a_{12}\\a_{21} & a_{22}\end{bmatrix} \in \cp{M}_{2 \times 2}(\real).
        \]
        Então:
        \begin{align*}
            (\alpha\beta)\cdot A & = (\alpha\beta)\begin{bmatrix} a_{11} & a_{12}\\a_{21} & a_{22}\end{bmatrix}
            \\ &= \begin{bmatrix} (\alpha\beta)a_{11} & (\alpha\beta)a_{12}\\(\alpha\beta)a_{21} & (\alpha\beta)a_{22}\end{bmatrix}
            \\ &= \begin{bmatrix} \alpha(\beta a_{11}) & \alpha(\beta a_{12})\\\alpha(\beta a_{21}) & \alpha(\beta a_{22})\end{bmatrix}
            \\ &= \alpha\cdot\begin{bmatrix} \beta a_{11} & \beta a_{12}\\\beta a_{21} & \beta a_{22}\end{bmatrix}
            \\ &= \alpha\cdot\left(\beta\cdot\begin{bmatrix} a_{11} & a_{12}\\a_{21} & a_{22}\end{bmatrix}\right)
            \\ &= \alpha\cdot(\beta\cdot A)
        \end{align*}

        \item[M2)] Seja $1 \in \real$ e
        \[
        A = \begin{bmatrix} a_{11} & a_{12}\\a_{21} & a_{22}\end{bmatrix} \in \cp{M}_{2 \times 2}(\real).
        \]
        Temos
        \begin{align*}
            1\cdot A & = \begin{bmatrix} a_{11} & a_{12}\\a_{21} & a_{22}\end{bmatrix}
            \\ &= \begin{bmatrix} 1\cdot a_{11} & 1\cdot a_{12}\\1\cdot a_{21} & 1\cdot a_{22}\end{bmatrix}
            \\ &= \begin{bmatrix} a_{11} & a_{12}\\a_{21} & a_{22}\end{bmatrix} = A.
        \end{align*}

        \item Sejam $\alpha \in \real$ e
        \[
        A = \begin{bmatrix} a_{11} & a_{12}\\a_{21} & a_{22}\end{bmatrix} \in \cp{M}_{2 \times 2}(\real),
        B = \begin{bmatrix} b_{11} & b_{12}\\b_{21} & b_{22}\end{bmatrix} \in \cp{M}_{2 \times 2}(\real).
        \]
        Temos
        \begin{align*}
            \alpha\cdot(A + B) & = \alpha\cdot\left(
            \begin{bmatrix} a_{11} & a_{12}\\a_{21} & a_{22}\end{bmatrix} +
            \begin{bmatrix} b_{11} & b_{12}\\b_{21} & b_{22}\end{bmatrix}\right)
            \\ &= \alpha\cdot\left(\begin{bmatrix} a_{11} + b_{11} & a_{12} + b_{12}\\a_{21} + b_{21} & a_{22} + b_{22}\end{bmatrix}\right)
            \\ &= \begin{bmatrix} \alpha(a_{11} + b_{11}) & \alpha(a_{12} + b_{12})\\\alpha(a_{21} + b_{21}) & \alpha(a_{22} + b_{22})\end{bmatrix}
            \\ &= \begin{bmatrix} \alpha a_{11} + \alpha b_{11} & \alpha a_{12} + \alpha b_{12}\\\alpha a_{21} + \alpha b_{21} & \alpha a_{22} + \alpha b_{22}\end{bmatrix}
            \\ &= \begin{bmatrix} \alpha a_{11} & \alpha a_{12}\\ \alpha a_{21} & \alpha a_{22}\end{bmatrix}
            + \begin{bmatrix} \alpha b_{11} & \alpha b_{12}\\ \alpha b_{21} & \alpha b_{22}\end{bmatrix}
            \\ &= \alpha\begin{bmatrix} a_{11} & a_{12}\\a_{21} & a_{22}\end{bmatrix} +
            \alpha\begin{bmatrix} b_{11} & b_{12}\\b_{21} & b_{22}\end{bmatrix}
            \\ &= \alpha\cdot A + \alpha\cdot B
        \end{align*}

        \item[D2)] Sejam $\alpha$, $\beta \in \real$ e
        \[
        A = \begin{bmatrix} a_{11} & a_{12}\\a_{21} & a_{22}\end{bmatrix} \in \cp{M}_{2 \times 2}(\real).
        \]
        Temos
        \begin{align*}
            (\alpha + \beta)\cdot A & = (\alpha + \beta)\cdot\begin{bmatrix} a_{11} & a_{12}\\a_{21} & a_{22}\end{bmatrix}
            \\ &= \begin{bmatrix} (\alpha + \beta)a_{11} & (\alpha + \beta)a_{12}\\(\alpha + \beta)a_{21} & (\alpha + \beta)a_{22}\end{bmatrix}
            \\ &= \begin{bmatrix} \alpha a_{11} + \beta a_{11} & \alpha a_{12} + \beta a_{12}\\\alpha a_{21} + \beta a_{21} & \alpha a_{22} + \beta a_{22}\end{bmatrix}
            \\ &= \begin{bmatrix} \alpha a_{11} & \alpha a_{12}\\ \alpha a_{21} & \alpha a_{22}\end{bmatrix}
            + \begin{bmatrix} \beta a_{11} & \beta a_{12}\\ \beta a_{21} & \beta a_{22}\end{bmatrix}
            \\ &= \alpha\cdot\begin{bmatrix} a_{11} & a_{12}\\a_{21} & a_{22}\end{bmatrix}  + \beta\cdot\begin{bmatrix} a_{11} & a_{12}\\a_{21} & a_{22}\end{bmatrix}
            \\ &= \alpha\cdot A + \beta\cdot A.
        \end{align*}
    \end{itemize}
    Assim $V = \cp{M}_{2 \times 2}(\real)$ é um espaço vetorial sobre $\real$.
\end{document}
